\chapter{Troubleshooting and questions}
\label{ch:troubleshooting}

Every reply of the bot says what happened and what to do next. This chapter collects the
situations where that isn't enough, or where the problem is outside the bot.

\section{The bot doesn't react at all}
\label{sec:no-reply}

\begin{itemize}
  \item \textbf{You have no role yet.} The bot silently ignores everyone who isn't a
        professor, TA or bot admin, so that students can't probe it. Ask a bot admin to add
        you as professor, or the course owner to add you as TA.
  \item \textbf{The command doesn't start with \cmd{!}}, or there is something before it.
        Commands must be the first thing in the message.
  \item \textbf{You edited the message.} The bot doesn't read edits. Send the corrected
        command as a new message.
  \item \textbf{The bot is down.} If even \cmd{!whoami} gets no \robot{} reaction after a
        minute, tell a bot admin.
\end{itemize}

\section{Replies take long}
\label{sec:slow}

Watch for the \robot{} reaction:

\begin{itemize}
  \item \textbf{The reaction appears quickly, the reply later:} the bot is working. Creating a
        course with many rooms takes a few seconds; bulk work runs as a job and reports when
        it's done.
  \item \textbf{The reaction itself comes late:} your message reached the bot late. That is a
        server or network issue, not your command. If it persists, tell a bot admin the
        time; the bot logs for every command how long the message took to arrive.
\end{itemize}

\section{Messages from the bot and what to do}

\begin{center}
\begin{tabularx}{\textwidth}{@{}>{\raggedright\arraybackslash}p{5.3cm}L@{}}
\toprule
\textbf{The bot says \dots} & \textbf{What to do} \\
\midrule
\enquote{Please send me commands in a direct chat \dots} &
  Use your one-to-one chat with \botid{}. In course rooms the bot doesn't take commands. \\
\enquote{You are not allowed to do that: only \dots} &
  The message names who may do it. Check your roles with \cmd{!whoami}; ask the course
  owner or a bot admin. \\
\enquote{There is no plan with the code \dots} &
  The plan expired (10 minutes), was replaced by a newer one, or someone else's code was
  used. Send the original command again. \\
\enquote{\dots{} send the corrected \dots{} here} &
  Your file was rejected for the reason given above it. Fix the file and send only the
  file again. \\
\enquote{I couldn't read the arguments} &
  Something is missing or a title with spaces isn't in quotes. The reply shows the right
  form; \cmd{!help <command>} has an example. \\
\enquote{The Matrix server refused a step \dots} &
  The server rejected an action, e.g.\ inviting a deactivated account. The reason is in
  the message; if it doesn't help, try again later or ask an admin. \\
\enquote{Something went wrong on my side} &
  An unexpected error. Parts of the command may have run already; check with
  \cmd{!course info}. Tell a bot admin the time of your command. \\
\enquote{\dots{} a server check failed when I started} &
  A server problem (e.g.\ professor bots misconfigured). Tell a bot admin. \\
\bottomrule
\end{tabularx}
\end{center}

\section{Questions from students}

\begin{description}[style=nextline,leftmargin=1.5em,font=\normalfont\bfseries]
  \item[\enquote{My link says it is invalid or expired.}]
    Links expire after 14 days and work only once. If the student has no account yet:
    \cmd{!signup renew <email>} gives a new link. If the link was used already, they have an
    account: they log in with the username shown on the link page and their password.
  \item[\enquote{I forgot my password.}]
    Only a bot admin can create a reset link (\cmd{!signup reset}). Point the student to the
    admins.
  \item[\enquote{I have an account but don't see the course.}]
    Existing accounts get an \emph{invitation} to the course space, which they have to accept
    in Element (look under \ui{Invites} or in the space list). If there is no invitation, they
    weren't on your list: \cmd{!signup new <email> <course>} invites them right away.
  \item[\enquote{I registered with my staff address, the list has my student address.}]
    Doesn't matter: \cmd{name@unibas.ch} and \cmd{name@stud.unibas.ch} are the same account.
  \item[\enquote{I can't read older messages.}]
    Encrypted messages need the keys from their other devices or their recovery key. After
    a password reset or on a new device without the recovery key, older messages stay
    unreadable. That can't be fixed afterwards.
  \item[\enquote{I can't write in Announcements.}]
    That's intended: only staff can post there. Questions belong in \enquote{Questions}.
\end{description}

\section{Things the bot doesn't do}

\begin{itemize}
  \item It doesn't send emails; you send the signup links (\autoref{sec:mailmerge}).
  \item It doesn't delete courses or rooms; \cmd{!room remove} only takes a room out of the
        course.
  \item It doesn't know when someone leaves the university. Remove students from a course
        with \cmd{!roster sync --remove}; deactivating accounts is up to the admins.
  \item It doesn't answer in course rooms, and it ignores students.
\end{itemize}
